% =====================================================================
\section{Последовательный порт: разговор с компьютером}
\label{les:serial}
% =====================================================================
\lessonintro
  {понять, как биты путешествуют по одному проводу, и научиться принимать команды с компьютера}
  {урок~\ref{les:signals} (биты и байты), урок~\ref{les:ide} (\code{Serial.print}), урок~\ref{les:pwm} (яркость)}
  {светодиод слушается команд \code{on}, \code{off} и \code{level 30}; игра «Угадай число»}
  {схема со светодиодом на \pin{4} (рис.~\ref{fig:led}); для последнего раздела --- модуль GPS или вторая плата (по желанию)}

С самого первого урока мы пользуемся Serial Monitor: плата присылает сообщения, и мы их читаем.
Пора разобраться, как это устроено, и научиться разговаривать в обратную сторону --- отправлять
плате команды.

\subsection{Как передать букву по одному проводу}

Внутри компьютера каждая буква --- это число. Например, латинская «A» --- это 65, а в двоичном виде
\code{01000001}. Эти 8 бит (1 байт) можно передать по одному проводу друг за другом, как в
уроке~\ref{les:signals}. Такой способ называют \textbf{последовательной передачей}, а устройство внутри
чипа, которое этим занимается, --- \textbf{UART} (англ. \emph{Universal Asynchronous Receiver-Transmitter},
универсальный асинхронный приёмопередатчик).

Слово «асинхронный» значит, что у передатчика и приёмника нет общего «метронома». Поэтому они заранее
договариваются о \textbf{скорости}: сколько бит в секунду передавать. Эту скорость называют
\textbf{бодами}; мы используем 115\,200 бод.

\begin{figure}[H]
\centering
\begin{tikzpicture}[x=1.05cm,y=1cm]
  \def\hi{1}
  \draw[very thick,accent] (0,\hi) -- (1,\hi) -- (1,0) -- (2,0) -- (2,\hi) -- (3,\hi) -- (3,0) -- (8,0) -- (8,\hi) -- (9,\hi) -- (9,0) -- (10,0) -- (10,\hi) -- (12,\hi);
  \foreach \x in {1,...,11} \draw[gray!50,dashed] (\x,-0.3) -- (\x,1.4);
  \foreach \x/\l in {1.5/Старт,2.5/D0,3.5/D1,4.5/D2,5.5/D3,6.5/D4,7.5/D5,8.5/D6,9.5/D7,10.5/Стоп}
    \node[font=\sffamily\scriptsize] at (\x,1.65) {\l};
  \foreach \x/\b in {2.5/1,3.5/0,4.5/0,5.5/0,6.5/0,7.5/0,8.5/1,9.5/0}
    \node[font=\ttfamily\small,text=esp] at (\x,-0.5) {\b};
  \node[font=\sffamily\scriptsize] at (0.5,1.65) {покой};
  \node[font=\sffamily\scriptsize] at (11.5,1.65) {покой};
  \draw[decorate,decoration={brace,mirror,amplitude=4pt}] (2,-0.8) -- (10,-0.8)
    node[midway,below=5pt,font=\sffamily\small]{8 бит данных, младший бит первым: \texttt{01000001} = 65 = «A»};
  \draw[<->] (1,-1.8) -- node[below,font=\scriptsize]{$t_\text{бита} = 1/115\,200~\text{с} \approx 8{,}7$~мкс} (2,-1.8);
\end{tikzpicture}
\caption{Передача буквы «A». В покое линия в состоянии 1. Старт-бит (0) говорит «внимание, начинаю!»,
стоп-бит (1) --- «закончил»}
\end{figure}

\begin{fact}
При скорости 115\,200 бод на один символ уходит 10 бит (8 бит данных + старт + стоп), значит, за
секунду можно передать 11\,520 символов --- примерно 5 страниц этой книги.
\end{fact}

\subsection{Serial и Serial Monitor}

В ESP32-S3 три UART. В Arduino им соответствуют объекты:

\begin{center}\small
\begin{tabular}{@{}llll@{}}
\toprule
Объект Arduino & Блок & Выводы & Для чего \\
\midrule
\code{Serial} & USB или UART0 & разъём USB / 43, 44 & связь с компьютером \\
\code{Serial1} & UART1 & любые, задаются в \code{begin()} & свободен \\
\code{Serial2} & UART2 & любые, задаются в \code{begin()} & свободен \\
\bottomrule
\end{tabular}
\end{center}

\begin{figure}[H]
\centering
\begin{tikzpicture}[node distance=10mm]
  \node[blk,fill=blkmem,minimum width=30mm,minimum height=12mm] (pc) {Компьютер\\ \scriptsize Serial Monitor};
  \node[blk,fill=blkcpu,minimum width=30mm,minimum height=12mm,right=40mm of pc] (esp) {ESP32-S3\\ \scriptsize \code{Serial}};
  \draw[arr,accent] ($(pc.east)+(0,0.25)$) -- node[above,font=\ttfamily\scriptsize]{"on\textbackslash n"} ($(esp.west)+(0,0.25)$);
  \draw[arr,okgreen] ($(esp.west)+(0,-0.25)$) -- node[below,font=\ttfamily\scriptsize]{"OK\textbackslash n"} ($(pc.east)+(0,-0.25)$);
  \node[font=\sffamily\scriptsize,below=1mm of pc] {поле ввода вверху окна};
  \node[font=\sffamily\scriptsize,below=1mm of esp] {\code{Serial.readStringUntil()}};
\end{tikzpicture}
\caption{Разговор в две стороны: команда от компьютера и ответ платы}
\end{figure}

Основные функции:
\begin{itemize}
  \item \code{Serial.begin(115200)} --- включить порт; скорость в Serial Monitor должна совпадать;
  \item \code{Serial.print()}, \code{println()}, \code{printf()} --- отправить текст;
  \item \code{Serial.available()} --- сколько байт пришло и ждёт чтения;
  \item \code{Serial.readStringUntil('\textbackslash n')} --- прочитать строку до символа «новая строка».
\end{itemize}

\begin{tip}
В Serial Monitor выбери окончание строки \textsf{New Line}: тогда каждая команда будет
заканчиваться символом \code{\textbackslash n}, и \code{readStringUntil()} сразу поймёт, где конец команды.
\end{tip}

\subsection{Практика: светодиод слушается команд}

Используем схему со светодиодом на \pin{4} (рис.~\ref{fig:led}) и ШИМ из урока~\ref{les:pwm}. Функция
\code{handleSerial()} ничего не ждёт: если команды нет, она сразу возвращается, и \code{loop()} может
заниматься другими делами.

\codecap{Управление светодиодом командами}
\codeinput{l09_commands}

Попробуй: \code{on}, \code{level 10}, \code{level 100}, \code{status}, \code{off}.

\subsection{Учимся программировать: строки и цепочка условий}

Команда из Serial Monitor приходит как \textbf{строка} --- переменная типа \code{String}. Со строками можно
делать много полезного:

\begin{table}[H]
\centering\small
\begin{tabular}{@{}lll@{}}
\toprule
Запись & Что делает & Пример для \code{cmd = "level 30"}\\
\midrule
\code{cmd == "on"} & сравнить строки & \code{false}\\
\code{cmd.startsWith("level ")} & начинается ли с этих букв & \code{true}\\
\code{cmd.substring(6)} & вырезать часть, начиная с символа № 6 & \code{"30"}\\
\code{cmd.substring(6).toInt()} & превратить текст в число & \code{30}\\
\code{cmd.length()} & длина строки & \code{8}\\
\code{cmd.trim()} & убрать пробелы и переводы строки по краям & ---\\
\code{"уровень " + String(30)} & склеить строки & \code{"уровень 30"}\\
\bottomrule
\end{tabular}
\caption{Полезные действия со строками \code{String}. Символы нумеруются с нуля, как ячейки массива}
\end{table}

Чтобы выбрать одно действие из многих, условия выстраивают в \textbf{цепочку} \code{if} --- \code{else if} ---
\code{else}. Они проверяются по очереди сверху вниз, и выполняется только \emph{первая} подходящая ветка.
Последний \code{else} ловит всё остальное --- например, опечатки.

\begin{figure}[H]
\centering
\begin{tikzpicture}[node distance=5mm and 12mm]
  \node[fcio] (in) {прочитать строку cmd};
  \node[fcdec,below=of in] (q1) {cmd == "on" ?};
  \node[fcproc,right=of q1] (a1) {ledOn = true};
  \node[fcdec,below=of q1] (q2) {cmd == "off" ?};
  \node[fcproc,right=of q2] (a2) {ledOn = false};
  \node[fcdec,below=of q2] (q3) {начинается\\ с "level " ?};
  \node[fcproc,right=of q3] (a3) {level = число};
  \node[fcproc,below=of q3] (a4) {напечатать\\ список команд};
  \draw[arr] (in) -- (q1);
  \foreach \q/\a in {q1/a1,q2/a2,q3/a3}{\draw[arr] (\q) -- node[above,font=\sffamily\scriptsize]{да} (\a);}
  \draw[arr] (q1) -- node[left,font=\sffamily\scriptsize]{нет} (q2);
  \draw[arr] (q2) -- node[left,font=\sffamily\scriptsize]{нет} (q3);
  \draw[arr] (q3) -- node[left,font=\sffamily\scriptsize]{нет} (a4);
\end{tikzpicture}
\caption{Цепочка условий в функции \code{handleSerial()}: проверяем команды по очереди}
\end{figure}

Обрати внимание на \code{return} в начале функции: \code{if (!Serial.available()) return;} --- «если данных
нет, сразу выйти из функции». Такой ранний выход избавляет от лишних фигурных скобок.

\subsection{Практика: игра «Угадай число»}

Теперь, когда мы умеем читать ввод с компьютера, сделаем игру. Плата загадывает случайное число от 1 до 100,
а ты угадываешь, вводя числа в поле Serial Monitor (окончание строки --- \textsf{New Line}). Плата подсказывает
«мало» или «много». Светодиод для игры не нужен.

\begin{figure}[H]
\centering
\begin{tikzpicture}[node distance=6mm]
  \node[fcterm] (s) {начало};
  \node[fcproc,below=of s] (new) {загадать число secret,\\ попытки = 0};
  \node[fcio,below=of new] (in) {ввести guess};
  \node[fcproc,below=of in] (t) {попытки = попытки + 1};
  \node[fcdec,below=of t] (q1) {guess < secret ?};
  \node[fcio,left=10mm of q1] (m1) {«мало!»};
  \node[fcdec,below=9mm of q1] (q2) {guess > secret ?};
  \node[fcio,left=10mm of q2] (m2) {«много!»};
  \node[fcio,right=10mm of q2] (w) {«Угадал!»};
  \draw[arr] (s) -- (new); \draw[arr] (new) -- (in); \draw[arr] (in) -- (t); \draw[arr] (t) -- (q1);
  \draw[arr] (q1) -- node[above,font=\sffamily\scriptsize]{да} (m1);
  \draw[arr] (q1) -- node[right,font=\sffamily\scriptsize]{нет} (q2);
  \draw[arr] (q2) -- node[above,font=\sffamily\scriptsize]{да} (m2);
  \draw[arr] (q2) -- node[above,font=\sffamily\scriptsize]{нет} (w);
  \draw[arr] (m1.west) -- ++(-5mm,0) |- (in.west);
  \draw[arr] (m2.west) -- ++(-10mm,0) |- (in.west);
  \draw[arr] (w.east) -- ++(5mm,0) |- (new.east);
\end{tikzpicture}
\caption{Блок-схема игры «Угадай число»}
\end{figure}

\codecap{Игра «Угадай число»}
\codeinput{l09_guess}

Здесь встретилось всё, что мы уже знаем: глобальные \textbf{переменные} \code{secret} и \code{tries},
\textbf{цепочка условий}, цикл \code{while}, который ждёт ввода, и \textbf{свои функции} \code{newGame()} и
\code{readNumber()}. Новые готовые функции: \code{random(1, 101)} даёт случайное число от 1 до 100, а
\code{Serial.parseInt()} читает из пришедших символов целое число.

\begin{fact}
Если угадывать с умом --- каждый раз называть середину оставшегося промежутка (50, потом 25 или 75\ldots), ---
любое число от 1 до 100 угадывается не больше чем за 7 попыток, ведь $2^7 = 128 > 100$. Точно так же работает
АЦП из урока~\ref{les:signals}: он «угадывает» напряжение за 12 шагов. Этот способ называется
\textbf{двоичным поиском}.
\end{fact}


\subsection{Второй UART: связь с другими устройствами}

На том же языке «разговаривают» GPS-приёмники, модули связи и другие микроконтроллеры. Благодаря
матрице GPIO (урок~\ref{les:pins}) UART1 можно вывести на любые выводы, например \pin{17} и \pin{18}.
Главное правило: \textbf{TX одного устройства соединяется с RX другого} (что один говорит --- другой
слушает), а земли объединяются.

\begin{twoview}{Подключение модуля GPS ко второму UART}{fig:uart}
\begin{tikzpicture}[bb]
  \bbBoard{24}
  \bbDevKit{29}{25}
  \bbModule{8}{7}{4}{7}{green!40!black}{GPS}{VCC,GND,TX,RX}
  \draw[wire=wblue]  (dk-17) -- (27.3,16) -- (27.3,5) -- (11,5);
  \draw[wire=wgreen] (dk-18) -- (27.9,15) -- (27.9,3) -- (10,3);
  \draw[wire=wblack] (9,4) -- (9,1);
  \draw[wire=wblack] (dk-GND) -- (28.6,4) -- (28.6,1) -- (24,1);
  \draw[wire=wred]   (dk-3V3) -- (28.3,25) -- (28.3,27.3) -- (-1.5,27.3) -- (-1.5,6) -- (8,6);
\end{tikzpicture}
\nextview
\begin{circuitikz}
  \node[draw,thick,fill=blkcpu,minimum width=2.6cm,minimum height=3.2cm,label={[font=\sffamily\bfseries]above:ESP32-S3}] (a) at (0,0) {};
  \node[draw,thick,fill=blkper,minimum width=2.6cm,minimum height=3.2cm,label={[font=\sffamily\bfseries]above:Модуль GPS}] (b) at (7,0) {};
  \draw ($(a.east)+(0,1.1)$) node[left,font=\scriptsize]{3V3} -- ($(b.west)+(0,1.1)$) node[right,font=\scriptsize]{VCC};
  \draw ($(a.east)+(0,0.4)$) node[left,font=\scriptsize]{TX \pin{17}} -- ++(2.5,0) -- ($(b.west)+(-1,-0.3)$) -- ($(b.west)+(0,-0.3)$) node[right,font=\scriptsize]{RX};
  \draw ($(a.east)+(0,-0.3)$) node[left,font=\scriptsize]{RX \pin{18}} -- ++(2.5,0) -- ($(b.west)+(-1,0.4)$) -- ($(b.west)+(0,0.4)$) node[right,font=\scriptsize]{TX};
  \draw ($(a.east)+(0,-1.1)$) node[left,font=\scriptsize]{GND} -- ($(b.west)+(0,-1.1)$) node[right,font=\scriptsize]{GND};
\end{circuitikz}
\end{twoview}

\codecap{Мост: всё, что пришло с модуля, пересылается на компьютер, и наоборот}
\codeinput{l09_bridge}

\begin{caution}
Некоторые модули (часть Arduino, старые GPS) работают с логикой 5~В. Их TX нельзя подключать к RX
ESP32-S3 напрямую --- помнишь «опасную зону» из урока~\ref{les:signals}? Нужен делитель напряжения
или преобразователь уровней.
\end{caution}

\begin{summary}
\begin{itemize}[leftmargin=*]
  \item UART передаёт байты по одному проводу: старт-бит, 8 бит данных, стоп-бит.
  \item Скорость (бод) у обеих сторон должна совпадать.
  \item Команды читаем, не останавливая программу: сначала \code{available()}, потом \code{readStringUntil()}.
  \item TX одного устройства --- к RX другого, земли --- вместе.
  \item Строки \code{String} можно сравнивать, резать и превращать в числа; цепочка \code{if}/\code{else if}/\code{else} выбирает одно действие из многих.
\end{itemize}
\end{summary}

\begin{quiz}
\begin{enumerate}
  \item Зачем нужен старт-бит?
  \item Что будет, если в Serial Monitor выбрать скорость 9600, а в программе --- 115\,200?
  \item Как соединить TX и RX двух плат?
\end{enumerate}
\end{quiz}

\begin{tasks}
\begin{enumerate}
  \item Добавь команду \code{blink N}: мигнуть светодиодом N раз (без \code{delay()}).
  \item Раз в 10 секунд выводи время работы платы, не мешая приёму команд.
  \item Соедини две платы по UART: команды, набранные на одной, пусть управляют светодиодом другой.
  \item Измени игру: пусть плата отгадывает число, загаданное тобой, двоичным поиском. Ты отвечаешь «+», «−» или «=».
\end{enumerate}
\end{tasks}
